% Generated from locale/tr/content/model-theory/models-of-arithmetic/non-standard-models.tex; do not hand-edit.


\olfileid[tr]{mod}{mar}{nst}
\section{Standart Olmayan Modeller}

\begin{explain}
$\Lang{L_A}$ için !!a{structure}, $\Struct N$ ile izomorfsa standart
deriz. !!a{structure}, $\Struct N$ ile izomorf değilse standart olmayan
diye adlandırılır.
\end{explain}

\begin{defn}
$\Lang{L_A}$ için !!{structure}~$\Struct M$, $\Struct N$ ile izomorf
değilse \emph{standart olmayan} bir yapıdır. Bazı $n\in\Nat$ için
$\Value{\num n}{M}$ değerine eşit olan $x\in\Domain M$ !!{element}s
$\Struct M$ yapısının \emph{standart sayıları}; böyle olmayanlar ise
\emph{standart olmayan sayıları} diye adlandırılır.
\end{defn}

\begin{explain}
\olref[stm]{prop:standard-domain} gereği $\Lang{L_A}$ için her standart
!!{structure} yalnızca standart !!{element}s içerir. Dolayısıyla
standart olmayan !!a{structure} en az bir standart olmayan eleman
içermelidir. Aslında standart olmayan !!a{element} bulunması,
!!{structure} ifadesinin standart olmadığını güvence altına alır.
\end{explain}

\begin{prop}
$\Lang{L_A}$ için !!a{structure}~$\Struct M$ standart olmayan bir sayı
içeriyorsa $\Struct M$ standart değildir.
\end{prop}

\begin{proof}
Tersini, yani $\Struct M$ standart olduğu halde standart olmayan bir
sayı~$x$ içerdiğini varsayalım. $g\colon\Nat\to\Domain M$ bir
izomorfizma olsun. $n$ üzerinde tümevarımla
$g(\Value{\num n}{N})=\Value{\num n}{M}$ olduğu kolayca görülür.
Başka bir deyişle $g$, $\Struct N$ yapısının standart sayılarını
$\Struct M$ yapısının standart sayılarına eşler. $\Struct M$ standart
olmayan bir sayı içeriyorsa $g$ !!{surjective} olamaz; bu varsayımla
çelişir.
\end{proof}

\begin{prob}
$\Th Q$ kuramının şu aksiyomları içerdiğini hatırlayın:
\begin{align*}
& \lforall[x][\lforall[y][(\eq[x'][y'] \lif \eq[x][y])]] \tag{$!Q_1$}\\
& \lforall[x][\eq/[\Obj 0][x']] \tag{$!Q_2$}\\
& \lforall[x][(\eq[x][\Obj 0] \lor \lexists[y][\eq[x][y']])] \tag{$!Q_3$}.
\end{align*}
Aşağıdaki koşulları sağlayan !!{structure}s~$\Struct{M_1}$,
$\Struct{M_2}$ ve $\Struct{M_3}$ verin:
\begin{enumerate}
\item $\Sat{M_1}{!Q_1}$, $\Sat{M_1}{!Q_2}$, $\Sat/{M_1}{!Q_3}$;
\item $\Sat{M_2}{!Q_1}$, $\Sat/{M_2}{!Q_2}$, $\Sat{M_2}{!Q_3}$; ve
\item $\Sat/{M_3}{!Q_1}$, $\Sat{M_3}{!Q_2}$, $\Sat{M_3}{!Q_3}$.
\end{enumerate}
Her biri için yalnızca $\Assign{\Obj0}{M_i}$ ve
$\Assign{\prime}{M_i}$ yorumlarını belirtmeniz yeterlidir.
\end{prob}

\begin{explain}
$\Lang{L_A}$ için standart olmayan !!{structure}s vermek oldukça
kolaydır. Örneğin !!{domain}~$\Int$ olan ve bütün mantıksal olmayan
simgeleri alışılmış biçimde yorumlayan yapıyı alın. Negatif sayılar
hiçbir $n$ için $\num n$ değerleri olmadığından bu yapı standart
değildir. Elbette $\Struct N$ ile aynı tümceleri doğru kılma anlamında
bir aritmetik \emph{modeli} olmayacaktır. Örneğin
$\lforall[x][\eq/[x'][\Obj0]]$ yanlıştır. Bununla birlikte tıkızlık
teoremini kullanarak standart olmayan aritmetik modellerinin varlığını
kolayca ispatlayabiliriz.
\end{explain}

\begin{prop}
$\Th{TA}=\Setabs{!A}{\Sat{N}{!A}}$, $\Struct N$ yapısının kuramı olsun.
$\Th{TA}$ kuramının !!a{enumerable} standart olmayan modeli vardır.
\end{prop}

\begin{proof}
$\Lang{L_A}$ dilini yeni bir !!{constant}~$c$ ile genişletin ve
\[
\Gamma=\Th{TA}\cup\{\eq/[c][\num0],\eq/[c][\num1],
\eq/[c][\num2],\dots\}
\]
!!{sentence}s kümesini ele alın. $\Gamma$'nın herhangi bir
$\Struct{M^c}$ modeli, bütün $n\in\Nat$ için
$x\neq\Value{\num n}{M}$ olan $x=\Assign{c}{M}$ !!a{element}
içerir; dolayısıyla $x$ standart değildir. Ayrıca
$\Th{TA}\subseteq\Gamma$ olduğundan açıkça
$\Sat{M^c}{\Th{TA}}$. Yalnız $c$ simgesini unutarak $\Struct{M^c}$
yapısını $\Lang{L_A}$ için !!a{structure}~$\Struct M$ biçimine
getirirsek evreni yine standart olmayan $x$ elemanını içerir ve
$\Sat{M}{\Th{TA}}$ geçerliliğini korur. Sonuncusu, $c$ simgesinin
$\Th{TA}$ içinde geçmemesiyle güvence altındadır. Dolayısıyla
$\Gamma$'nın bir modeli olduğunu göstermek yeterlidir.

$\Gamma$'nın bir modeli bulunduğunu göstermek için tıkızlık teoremini
kullanırız. $\Gamma$'nın her sonlu alt kümesi sağlanabilirse $\Gamma$
da sağlanabilir. Gelişigüzel sonlu bir $\Gamma_0\subseteq\Gamma$
alalım. $\Gamma_0$, $\Th{TA}$ içinden bazı !!{sentence}s ve
$\eq/[c][\num n]$ biçiminde sonlu sayıda tümce içerir. Bu ikinci türden
hiç tümce yoksa $k=-1$; varsa
$\eq/[c][\num k]\in\Gamma_0$ olacak en büyük $k$ olsun.
$\Struct N$ yapısını $\Assign{c}{N_k}=k+1$ yorumunu ekleyerek
$\Struct{N_k}$ biçiminde genişletin. $\Sat{N_k}{\Gamma_0}$: Eğer
$!A\in\Th{TA}$ ise $\Struct{N_k}$, $c$ dışında her bakımdan
$\Struct N$ ile aynı ve $c$, $!A$ içinde geçmediğinden
$\Sat{N_k}{!A}$. Ayrıca $n\leq k$ ve
$\Value c{N_k}=k+1$ olduğundan
$\Sat{N_k}{\eq/[c][\num n]}$. Böylece $\Gamma$'nın her sonlu alt
kümesi sağlanabilirdir.
\end{proof}

